\chapter{General discussion and conclusions}
\label{ch:discussion}

\begin{instructions}
About 5 pages; Conclusion is 15\,\% of the mark. The marksheet asks for three things here:
\begin{itemize}
  \item insight into the significance of the results;
  \item a critical review of your performance against the stated plan, including goals and schedule
        (Section~\ref{sec:review});
  \item a clear summary of the outcomes, in context of the thesis definitions and the literature, with recommendations
        for continuing the work.
\end{itemize}
Write this chapter last but one, before the abstract.
\end{instructions}

\section{Summary}

\begin{instructions}
One paragraph per research question, each stating the finding, not the method, with its headline number:
\begin{itemize}
  \item RQ1: detector agreement predicts human acceptance on every recording (0--5\,\% $\rightarrow$ 40--63\,\%).
        Choosy detectors hold across recordings; loose ones measure the spectrum. Against 20 surrogate draws, drop v1
        and the GADF CNN beat every draw on 16 of 16 channels.
  \item RQ2 (queue 35, judged blind): every seed beats random windows of its own length, Fisher $p<0.0001$ each (for
        example drop30 100\,\% against 22\,\%). All 51 earlier yeses held when judged blind (95\,\% interval
        93--100\,\%).
  \item RQ3:
        \begin{itemize}
          \item MiniRocket on the tool's verdicts reaches AUC 0.971 leave-one-out, against a shuffle null of 0.49 at
                most.
          \item On the held-out M4 exam (queue 36, scored once), R4 reaches AUC 0.98 and balanced accuracy 0.93,
                against 0.59--0.66 AUC for the reference CNNs.
          \item A CNN trained on R4's own rows and folds matches R4 (pooled AUC 0.970--0.973 against 0.971), so the
                reference CNNs' gap is their labels.
          \item The \emph{S.\ commune} exam waits on queue 37.
          \item Cluster labels do not replace human labels.
        \end{itemize}
\end{itemize}
Then one paragraph each for Chapter~\ref{ch:events} (for example: 91\,\% of accepted events are falls by the
fastest-edge rule) and Chapter~\ref{ch:roundb} (from Table~\ref{tab:rbsummary}).
\end{instructions}

\section{Discussion}

\begin{instructions}
Points to link across the chapters:
\begin{itemize}
  \item Do two independent routes to ``these belong together'' agree? Compare shape families (Library) with seed-search
        matches, and with the classifier's high-confidence events.
  \item Does what detectors miss explain what the classifiers get wrong? Compare R4's errors with the low-agreement
        strata and the miss-rate re-look list.
  \item Which techniques were useful, and which were dead ends:
        \begin{itemize}
          \item useful: the choosy drop detectors, detector agreement, seed search, MiniRocket, AAFT nulls;
          \item dead ends: the whole-signal matrix-profile threshold (0\,\% precision); 1-minute CNNs (chance); the
                band chain against its null; the fusion encoding not beating its components; the early
                binned-difference forecasting attempt (report it with its baseline);
          \item the Dehshibi detector: faithful to its authors' code but unselective on these recordings;
          \item \texttt{cwt\_events} against AAFT (B9): it finds no more structure than its input's values give.
        \end{itemize}
  \item \textbf{Modelling, only if it is of interest} (parked from Chapter~\ref{ch:roundb}, 4 Nov). A fitted model is
        not evidence on these signals:
        \begin{itemize}
          \item \textbf{They fit noise too.} FitzHugh--Nagumo, Izhikevich, van der Pol and Morris--Lecar all fit the
                mean drop cycle (R$^2$ mostly above 0.9), and fit fake cycles cut from noise about as well.
          \item \textbf{They do not forecast.} None predicts a held-out cycle better than the mean of its neighbours,
                on any of 11 trains of up to 110 cycles.
          \item \textbf{Their PRCs say nothing about the fungus.} The models' phase response curves are mostly type
                II, but 8 of 33 flip type when a different hidden coordinate is held fixed \cite{wilshin2024prc}.
          \item \textbf{Sources:} \src{docs/thesis/parked/ch10_b3_models.tex} and
                \src{docs/rq_roundA_2/findings/b3-long-trains-2026-11-03.md}.
        \end{itemize}
  \item What ``normalisation is not invariance'' means for any shape vocabulary.
  \item Set the results against the literature: no fungal study reports precision against blind human judgement, or
        against a surrogate.
\end{itemize}
\end{instructions}

\section{Critical review against the plan}
\label{sec:review}

\begin{table}[htbp]
\caption{The thesis plan (August 2026) against what was done.}
\label{tab:planreview}
\footnotesize
\begin{tabularx}{\textwidth}{p{3.8cm} p{3.6cm} X}
\toprule
\textbf{Planned} & \textbf{Outcome} & \textbf{Reason}\\
\midrule
Six RQs in three studies & Three RQs answered; three plan RQs partly or as pilots & Comparison at scale subsumed the chain-against-chain questions; scale and cross-channel questions need judging not yet done\\
\addlinespace
Feature freeze, 28 Aug & Not met; built to 21 Oct & The Panel interface did not scale and was rebuilt in React (ADR 0001); the deadline moved on 23 Sep\\
\addlinespace
Held-out M4, untouched until a single final run & Exposed to detector output 17--21 Oct; one-shot exam frozen 26 Oct & The lock was switched off for the demonstrations; disclosed in Table~\ref{tab:exposure}\\
\addlinespace
A surrogate on every result & On every reported result; one draw in the sweep & Cost; 20-draw re-runs prepared\\
\addlinespace
Morphology vocabulary tested (RQ5) & Not testable & No verdict carries a shape tag\\
\addlinespace
Library similarity edges persisted & Schema only, 0 edges & Families persisted through groupings instead\\
\addlinespace
Block UI generated from HoloViz Param & From \texttt{ParamSpec} in React & Panel retired\\
\addlinespace
Cluster labels as the route to training data & Human verdicts in blind queues & Cluster labels mapped to interesting/not gave many false alarms (Chapter~\ref{ch:rq3})\\
\bottomrule
\end{tabularx}
\end{table}

\begin{table}[htbp]
\caption{The plan's schedule against what happened.}
\label{tab:schedulereview}
\footnotesize
\begin{tabularx}{\textwidth}{p{2.6cm} p{4.4cm} X}
\toprule
\textbf{Period} & \textbf{Planned} & \textbf{What happened}\\
\midrule
To 28 Aug & Feature freeze & Orchestrated build of the first, Panel-based version (Claude, then DeepSeek tickets)\\
Late Aug -- mid Sep & Studies 1 and 2 on the primary recording & Drop-motif study (25 Aug -- 4 Sep): cross-species scale collapse and claims ledger; React stack chosen (15 Sep)\\
Mid Sep -- early Oct & Study 3 adjudication & Panel retired (21 Sep); data loss and recovery (21 Sep); round-A RQ files (from 3 Oct); RQ1 baseline (5 Oct)\\
Early Oct & M4 end-to-end run & Lock enforced; detector sweep and blind queues on five recordings (17--20 Oct); R4 (20 Oct); demonstration and poster (21 Oct)\\
Oct -- Nov & Results and discussion written & Blind seed panel and M4 exam (26 Oct on); round B (29 Oct -- 3 Nov); writing (4--8 Nov)\\
November & Draft, revision, submission & Draft 5 Nov; submitted by 9 Nov\\
\bottomrule
\end{tabularx}
\end{table}

\begin{instructions}
Write the review in your own voice, two or three paragraphs. Be honest about the main trade: much of the year went into
the instrument. Say what that bought: comparable runs, blind queues, nulls and recipes, which made the later results
quick to obtain (the sweep and R4 took days, not weeks). Say what it cost: several plan questions are pilots. Name what
you would do differently: freeze the tool earlier, keep the held-out lock on, and record CNN training provenance.
\end{instructions}

\section{Limitations}

\begin{instructions}
\begin{itemize}
  \item A single labeller (intra-rater $\kappa$ 0.84); no second rater.
  \item Mostly one species and one lab, with cross-species data limited to the drop-motif study and public recordings.
  \item Detector-proposed candidates only: there is no free-scan recall.
  \item One-draw surrogates in the sweep (20 draws on M2\_aug only).
  \item The Ward groupings are fragile.
  \item The M4 exposure.
  \item Shape normalisation does not remove resolution.
  \item Environment covariates on one recording, with unverified alignment.
\end{itemize}
\end{instructions}

\section{Future work}
\label{sec:futurework}

\begin{instructions}
Carried from the plan and updated:
\begin{itemize}
  \item \textbf{Forecasting}, with a baseline. The binned next-point attempt was highly susceptible to noise; report
        most-frequent-bin accuracy as its baseline. Forecasting is the precondition for interpretation.
  \item \textbf{Interpretation}, in that order: only a structure that predicts what follows is worth decoding.
  \item \textbf{Cross-species validation}, now possible: a second unseen-recording exam on \emph{S.\ commune}, and seeds
        across recordings.
  \item \textbf{Multivariate analysis and environment}: an Explore environment view (about 2--2.5 days as costed),
        covariate series in registration, and event rate against temperature.
  \item \textbf{Perturbation experiments}: measured phase response curves would need stimulation \cite{wilshin2024prc}.
  \item \textbf{A shape-class model}: a classifier over verified named Library classes with an explicit background
        class, needing 100--200 verified members per class.
  \item \textbf{Library mechanics}: event extents and merging partial motifs; an average-linkage option; persisted
        similarity edges.
  \item \textbf{Statistics}: 20-draw nulls beyond M2\_aug, whose 16 channels are the only ones that have them
        (Chapter~\ref{ch:rq1}); B1's interval read on every recording, with Benjamini--Hochberg across channels.
  \item \textbf{Round B's open questions}, unless they return before submission: B7 (multiscale entropy), B8 (motif
        prevalence) and B11 (the newer datasets); the \emph{S.\ commune} exam.
  \item \textbf{Method extensions}: the Morse-wavelet component of the Dehshibi detector as a reusable feature;
        interval temporal logic for sequences of windows; diffusion models for synthetic training data.
\end{itemize}
\end{instructions}

\section{Conclusions}

\begin{instructions}
Write last. Half a page. Answer the broad question in one sentence: recurring structure can be detected and
characterised, with stated precision against human judgement and a stated null. Then name the three things a reader
should take away, and the one thing the field should do next: report detector precision against blind judgement and
surrogates as standard.
\end{instructions}
